% Exact quintic case; independently audited by parent /root.
\subsection{The exact group for ordinary plane quintics in characteristic two}
\label{subsec:quintic-exact}

\begin{theorem}\label{thm:quintic-char2-monodromy}
For the universal ordinary smooth plane quintic in characteristic two,
the geometric and arithmetic monodromy groups on
$V=H^1_{\mathrm{et}}(C_{\bar\eta},\mathbf F_2)$ are both the full
stabilizer of the canonical nonzero covector:
\[
 G_{\mathrm{geom}}=G_{\mathrm{arith}}
 =\operatorname{Stab}_{\mathrm{GL}_6(\mathbf F_2)}(\ell)
 \simeq\mathbf F_2^5\rtimes\mathrm{GL}_5(\mathbf F_2).
\]
In particular their order is $319979520$, the representation is
reducible with invariant hyperplane of dimension five, and the
determinant character is trivial.
\end{theorem}
\begin{proof}

\paragraph{Theta-characteristic input, with a proof.}
For a smooth curve over an algebraically closed field of characteristic
two, let $B=F_*\mathcal O_C/\mathcal O_C$ for absolute Frobenius.
This is a line bundle with $B^{\otimes2}\simeq\omega_C$.
Locally a separating parameter $t$ gives the basis $1,t$ of
$F_*\mathcal O_C$ over $\mathcal O_C$; a change of separating parameter
$u=a^2+tb^2$ sends the generator of $B$ to $b$ times that generator,
whereas $du=b^2dt$.  This proves the claimed square relation.
The long exact sequence of
$0\to\mathcal O_C\to F_*\mathcal O_C\to B\to0$ shows
\[
 H^0(C,B)=\ker\bigl(F:H^1(C,\mathcal O_C)\to H^1(C,\mathcal O_C)\bigr).
\]
Thus $B$ is non-effective for an ordinary curve.
Every other theta characteristic is $B\otimes T$ for a nontrivial
two-torsion line bundle $T$.  Tensoring this sequence by $T$ and using
the projection formula and $F^*T=T^{\otimes2}\simeq\mathcal O_C$
gives an injection
\[
 H^0(C,\mathcal O_C)=H^0(C,F_*F^*T)
 \lhook\joinrel\longrightarrow H^0(C,B\otimes T),
\]
because $H^0(C,T)=0$.  Hence every theta characteristic other than $B$
is effective.  This is the statement recorded as St\"ohr--Voloch's
theorem in Ishitsuka--Ito, Theorem~3.1; Mumford (1971), p.~191,
describes the canonical characteristic-two theta characteristic.
For the geometric generic curve, the same reasoning may be performed
after perfecting its coefficient field.  Perfection does not change
the absolute Galois group, and $B$ is functorial under that group.

An ordinary curve of genus six has exactly $2^6=64$ geometric theta
characteristics.  On a plane quintic, $L=\mathcal O_C(1)$ is one of
them and $h^0(L)=3$, while $B$ is the unique non-effective one.
Under the identification of their set with the geometric points of
$J[2]$ by $T\mapsto L\otimes T$, the two distinguished points are
$0$ and $v=B\otimes L^{-1}\ne0$.  This is the dual representation
$V^\vee$, and $v$ is fixed by monodromy.

\paragraph{An irreducible incidence correspondence of contact conics.}
Let $S=\mathbf P(H^0(\mathbf P^2,\mathcal O(5)))$, of dimension twenty.
Let $I$ parametrize pairs $(C,Q)$ where $Q$ is a smooth plane conic
and the divisor cut out by $C$ on $Q$ is twice a divisor of degree five.
For a fixed conic $Q\simeq\mathbf P^1$, this means that the restriction
of its quintic equation, a binary form of degree ten, is a square.
The square forms constitute a vector subspace of dimension six inside
$H^0(Q,\mathcal O_Q(5))$, of dimension eleven.  The restriction map
from ternary quintics is surjective, with kernel
$Q\cdot H^0(\mathbf P^2,\mathcal O(3))$ of dimension ten.
Consequently the allowed quintic equations form a vector space of
dimension sixteen, or a projective space of dimension fifteen.
The smooth-conic parameter space has dimension five, so $\dim I=20$.
It is irreducible: fix $Q_0=(XZ+Y^2=0)$, take the irreducible product
$\mathrm{PGL}_3\times\mathbf P^{15}$ for its allowed quintics, and
map it surjectively to $I$ by changing projective coordinates.

The locus of pairs with a reduced reducible conic has dimension at
most nineteen.  Indeed, write $Q=L_1\cup L_2$ and $P=L_1\cap L_2$.
If $C$ is smooth and $C\cap Q$ is even, then $P\in C$: away from $P$
each of the two degree-five line sections must have even multiplicity,
so its multiplicity at $P$ is odd.  On each line the restriction is
therefore its linear equation for $P$ times the square of a binary
quadratic.  Each such restriction space has dimension three, and the
two restrictions automatically agree with value zero at $P$.
The combined restriction space has dimension six, so the quintic
equations again form a sixteen-dimensional vector space.
Pairs of distinct lines have dimension four, proving the bound
$4+15=19$.  Double-line conics cut exactly the divisors corresponding
to $L=\mathcal O_C(1)$.

Every effective theta characteristic $\theta$ and nonzero section $s$
gives a contact conic: $s^2\in H^0(C,\omega_C)=H^0(C,\mathcal O_C(2))$
lifts uniquely to a quadratic form on $\mathbf P^2$.
If that conic is a double line, then $\theta=L$.
Outside the dimension-at-most-nineteen image just described, each of
the other sixty-two theta characteristics therefore supplies a smooth
contact conic.  Thus $I\to S$ is dominant.  Since both dimensions are
twenty, its geometric generic fibre is zero dimensional.
It follows also that every such theta characteristic has $h^0=1$:
otherwise its projective space of sections would give a positive
dimensional family of conics in the same fibre.  Distinct theta
characteristics give distinct conics, because the half of an even
divisor on a smooth curve is unique.  The geometric generic fibre
therefore has exactly sixty-two points.

Irreducibility of $I$ implies transitivity of the geometric Galois
group on these sixty-two points.  This remains valid if the incidence
map has a purely inseparable part: the number of geometric generic
points is its separable degree, and a field's embeddings into a
separable closure form one orbit.  The Galois-equivariant
theta-characteristic interpretation proves that
$G_{\mathrm{geom}}$ on $V^\vee\setminus\{0,v\}$ is transitive.

\paragraph{The factor seven supplied by a quartic boundary.}
Take an ordinary smooth plane quartic $C_4$ and a line meeting it
transversely in four distinct points.  Their union $C_0$ is an ordinary
nodal quintic: the normalization sequence gives bijective Frobenius
on its $H^1(\mathcal O)$, of dimension six.  Over the parameter space
of such pairs, the etale normalization sequence gives a surjection
of local systems
\[
 H^1_{\mathrm{et}}(C_0,\mathbf F_2)
 \longrightarrow H^1_{\mathrm{et}}(C_4,\mathbf F_2).
\]
Forgetting the line is a smooth surjection with geometrically
connected fibres (nonempty opens of the dual plane), hence is
surjective on fundamental groups.  The quartic local system has
geometric monodromy $\mathrm{GL}_3(\mathbf F_2)$ by the already proved
quartic case.  Its order is divisible by seven, so the monodromy on
the nodal quintic family, and consequently the generic quintic group,
has order divisible by seven.  This specialization is legitimate
because the entire rank-six Lang system extends across $C_0$ on the
invertible-Hasse--Witt locus $W_5$.  The dense smooth open and $W_5$
have the same monodromy image for that extended finite etale system.

\paragraph{Finite group conclusion.}
Write $M=\operatorname{Stab}(v)\simeq2^5\!:\!\mathrm{GL}_5(2)$
in its action on $V^\vee$, and let $G\subset M$ be the geometric image.
Its transitive orbit of size sixty-two implies $31\mid|G|$; the
boundary argument gives $7\mid|G|$.
The complete maximal-subgroup list for $M$ is given in
Basheer--Moori (2021), Table~1, p.~3:
\[
\begin{array}{c|r|r}
\text{maximal subgroup}&\text{order}&\text{index in }M\\\hline
2_+^{1+8}:\mathrm{GL}_4(2)&10321920&31\\
2^{4+5}:\mathrm{GL}_4(2)&10321920&31\\
\mathrm{GL}_5(2)&9999360&32\\
2^{2+9}:(\mathrm{GL}_3(2)\times S_3)&2064384&155\\
2^9:(\mathrm{GL}_3(2)\times S_4)&2064384&155\\
(2^5:31):5&4960&64512.
\end{array}
\]
The first, second, fourth, and fifth orders are not divisible by
thirty-one.  The last is not divisible by seven.  The remaining
maximal subgroup is a complement $\mathrm{GL}_5(2)$; its action on
$V^\vee\setminus\{0,v\}$ has two orbits of length thirty-one, so it
cannot contain a group transitive on all sixty-two points.
No proper maximal subgroup can contain $G$, and therefore $G=M$.
The arithmetic image is contained in $M$ and contains $G$, hence
equals it.
\end{proof}

\paragraph{Independent elementary check of the complement issue.}
The same conclusion can be obtained from the maximal-subgroup table
for $\mathrm{GL}_5(2)$ in the ATLAS and the elementary fact
$H^1(\mathrm{GL}_m(2),\mathbf F_2^m)=0$ for $m\geq4$.
For completeness, let $c$ be a cocycle and put
$t_{ij}=1+E_{ij}$.  The relation $t_{ij}^2=1$ says that the $j$th
coordinate of $c(t_{ij})$ is zero.  Commuting $t_{ij}$ with
$t_{k\ell}$ for an auxiliary index $k$ distinct from $i,j,\ell$
shows that every coordinate except the $i$th is zero.
Thus $c(t_{ij})=a_{ij}e_i$.  The relation
$t_{ij}t_{jk}t_{ij}=t_{jk}t_{ik}$, for distinct indices,
gives $a_{ik}=a_{jk}$, so $a_{ij}=b_j$ depends only on $j$.
The coboundary associated with $b=(b_j)$ agrees with $c$ on all
elementary transvections, which generate the group.  Hence $c$ is
a coboundary.  The same result holds for the dual module by the
inverse-transpose automorphism.  Complements in $2^5:\mathrm{GL}_5(2)$
are therefore conjugate to the standard one, justifying the asserted
two-orbit description independently of the affine maximal-subgroup
classification.
